% ============================================================
% Section 2: Theoretical Framework
% ============================================================

\section{Theoretical Framework}
\label{sec:theory}

\subsection{Helmholtz Equation for the Optical Field}
\label{subsec:helmholtz}

The propagation of coherent light in a homogeneous and isotropic medium is governed
by the scalar Helmholtz equation:
%
\begin{equation}
  \nabla^2 E(\mathbf{r}) + k^2 E(\mathbf{r}) = 0,
  \label{eq:helmholtz}
\end{equation}
%
where $E(\mathbf{r}) \in \mathbb{C}$ is the complex amplitude of the electric field,
$k = 2\pi/\lambda$ is the wavenumber and $\lambda$ the laser wavelength.
In two spatial dimensions, with $\mathbf{r} = (x, y)$, the equation takes the form:
%
\begin{equation}
  \frac{\partial^2 E}{\partial x^2} + \frac{\partial^2 E}{\partial y^2} + k^2 E = 0.
  \label{eq:helmholtz2d}
\end{equation}

Since neural networks operate on real numbers, the complex field is decomposed into
its real and imaginary parts: $E = E_{\mathrm{real}} + i\,E_{\mathrm{imag}}$.
Both components satisfy Equation~\eqref{eq:helmholtz2d} independently:
%
\begin{equation}
  \nabla^2 E_{\mathrm{real}} + k^2 E_{\mathrm{real}} = 0, \qquad
  \nabla^2 E_{\mathrm{imag}} + k^2 E_{\mathrm{imag}} = 0.
  \label{eq:helmholtz_split}
\end{equation}
%
The speckle intensity, which is the physically observable quantity, is computed as:
%
\begin{equation}
  I(x,y) = |E|^2 = E_{\mathrm{real}}^2 + E_{\mathrm{imag}}^2.
  \label{eq:intensity}
\end{equation}

\subsection{Physics-Informed Neural Networks (PINNs)}
\label{subsec:pinns}

PINNs, introduced by Raissi et al. \cite{Raissi2019_PINNs} and extensively reviewed
by Karniadakis et al. \cite{Karniadakis2021_review} and Cuomo et al.
\cite{Cuomo2022_review}, approximate the solution of a PDE by means of a neural
network $\mathcal{N}_\theta(\mathbf{r})$ whose parameters $\theta$ are optimized by
minimizing a composite loss function:
%
\begin{equation}
  \mathcal{L}(\theta) = \mathcal{L}_{\mathrm{data}}(\theta)
                      + \lambda\,\mathcal{L}_{\mathrm{physics}}(\theta),
  \label{eq:loss_total}
\end{equation}
%
where $\mathcal{L}_{\mathrm{data}}$ penalizes the error on the boundary conditions
and $\mathcal{L}_{\mathrm{physics}}$ penalizes the violation of the PDE inside the
domain.
For the 2D Helmholtz equation with a complex field:
%
\begin{align}
  \mathcal{L}_{\mathrm{data}}(\theta) &=
    \frac{1}{N_b}\sum_{j=1}^{N_b}
    \left\|\mathcal{N}_\theta(\mathbf{r}_j^b) - E(\mathbf{r}_j^b)\right\|^2,
  \label{eq:loss_data} \\
  \mathcal{L}_{\mathrm{physics}}(\theta) &=
    \frac{1}{N_c}\sum_{i=1}^{N_c}
    \left[R_{\mathrm{real}}(\mathbf{r}_i^c)^2 + R_{\mathrm{imag}}(\mathbf{r}_i^c)^2\right],
  \label{eq:loss_physics}
\end{align}
%
where $R_{\mathrm{real}} = \nabla^2 E_{\mathrm{real}} + k^2 E_{\mathrm{real}}$ and
$R_{\mathrm{imag}} = \nabla^2 E_{\mathrm{imag}} + k^2 E_{\mathrm{imag}}$ are the
Helmholtz residuals, $\{\mathbf{r}_j^b\}_{j=1}^{N_b}$ are the boundary points and
$\{\mathbf{r}_i^c\}_{i=1}^{N_c}$ are the interior collocation points.
The derivatives $\nabla^2 E$ are obtained through automatic differentiation
\cite{Baydin2018_autodiff}.

\subsection{SIREN Networks}
\label{subsec:siren}

Standard activations (ReLU, tanh) produce unstable second-order derivatives and
exhibit \textit{spectral bias}: they tend to learn low frequencies before high ones,
which hinders the capture of the fine details of the speckle pattern.
Sitzmann et al. \cite{Sitzmann2020_SIREN} propose SIREN networks (Sinusoidal
Representation Networks), which use the activation $\phi(z) = \sin(\omega_0 z)$
together with a specific initialization that preserves the distribution of
activations across all layers.

The SIREN architecture with $L$ hidden layers of dimension $d$ is defined as:
%
\begin{equation}
\label{eq:siren}
\begin{aligned}
  \mathbf{h}_0 &= \mathbf{x}, \\
  \mathbf{h}_\ell &= \sin\!\left(\omega_0\,(\mathbf{W}_\ell\,\mathbf{h}_{\ell-1} + \mathbf{b}_\ell)\right),
  \quad \ell = 1, \ldots, L
\end{aligned}
\end{equation}
%
where $\omega_0$ is the frequency of the input layer and the matrices
$\mathbf{W}_\ell$ are initialized according to
$\mathcal{U}(-\sqrt{6/d},\,\sqrt{6/d})$ for the intermediate layers and
$\mathcal{U}(-1/n_{\mathrm{in}},\,1/n_{\mathrm{in}})$ for the first layer
\cite{Sitzmann2020_SIREN}.
This initialization ensures that the activation distributions are approximately
uniform in $[-1, 1]$, a necessary condition for the stability of derivatives of
arbitrary order.

\subsection{Statistics of Fully Developed Speckle}
\label{subsec:speckle}

When a rough surface with phase variation $\Delta\phi \gg 2\pi$ is illuminated by
coherent light, the contributions of the different points of the surface interfere
constructively and destructively in a random manner.
By the central limit theorem, the resulting complex field
$E = E_{\mathrm{real}} + i\,E_{\mathrm{imag}}$ follows a circular complex Gaussian
distribution \cite{Goodman2007_Speckle}:
%
\begin{equation}
  E_{\mathrm{real}},\; E_{\mathrm{imag}} \sim \mathcal{N}(0, \sigma_E^2),
  \quad \text{independent.}
  \label{eq:gaussian_field}
\end{equation}

The intensity $I = |E|^2$ then follows a negative exponential distribution:
%
\begin{equation}
  p(I) = \frac{1}{\langle I \rangle}\exp\!\left(-\frac{I}{\langle I \rangle}\right),
  \quad I \geq 0,
  \label{eq:pdf_intensity}
\end{equation}
%
with expected value $\langle I \rangle = 2\sigma_E^2$ and standard deviation
$\sigma_I = \langle I \rangle$.
The \textit{speckle contrast} is:
%
\begin{equation}
  C = \frac{\sigma_I}{\langle I \rangle} = 1,
  \label{eq:contrast}
\end{equation}
%
which equals exactly 1 for fully developed speckle and decreases when the field
loses coherence or when the surface roughness is insufficient.
This statistical criterion will be the main validation metric for speckle
generation in future work, complemented by the Kolmogorov-Smirnov (KS) test against
the theoretical exponential distribution.
